\documentclass[11pt,compress,t,notes=noshow, xcolor=table]{beamer}

\input{../../style/preamble}
\input{../../latex-math/basic-math}
\input{../../latex-math/basic-ml}
\input{../../latex-math/optim-basics}

% Lipschitz constant of the Hessian (Boyd's L; plain L is already the smoothness
% constant in this lecture, see 01-foundations/03-convex.tex)
\newcommand{\LH}{L_{\H}}

\title{Optimization in Machine Learning}

\begin{document}

\titlemeta{
Second order methods
}{
Deep Dive: Convergence of Newton-Raphson
}{
figure/NR_phases.png
}{
\item Assumptions
\item Damped phase
\item Quadratically convergent phase
}

\begin{framei}{Setting}
\item \textbf{Recall:}
\begin{itemizeM}
\item with full steps ($\alpha^{[t]} = 1$), NR may overshoot/diverge far from
$\xvs$ $\Rightarrow$ damping is necessary
\item many objective functions are only locally quadratic
\end{itemizeM}
\item In this deep dive, we show that damped NR has two phases:
\begin{itemizeM}
\item \textbf{Damped phase:} quadratic model unreliable, $\alpha^{[t]} < 1$ possible; $f$ decreases by a
\textit{constant} $\delta$ per step
\item \textbf{Pure Newton phase:} quadratic model accurate, $\alpha^{[t]} = 1$ always accepted; \textit{quadratic} convergence
\end{itemizeM}
\splitV[0.62]{
\imageC[0.85]{figure/NR_phases.png}
}{
{\footnotesize \textbf{Example 2} (cont'd):
$$f(\xv) = (1 + x_1^2)^{1/2} + \frac{x_2^2}{2}$$
from $\xv^{[0]} = (3, 2)^\top$,\\
with $\gamma_1 = 0.1$, $\tau = 0.7$}
}
\end{framei}


\begin{framei}[fs=small]{Setting}
\item \textbf{Assumptions}:
\begin{itemizeM}
\item $f \in \CC{2}$, $m$-strongly convex ($\H(\xv) \succeq m \id$) and $L$-smooth ($\H(\xv) \preceq L \id$)
\item $\H$ is \textbf{Lipschitz} with constant $\LH$ ($\rightarrow$ bound on the \textit{3rd} derivative):
$$\|\H(\xv) - \H(\xtil)\|_2 \le \LH \|\xv - \xtil\| \qquad \text{for all } \xv, \xtil$$
Interpretation: $\LH$ measures well the quadratic model approximates $f$ ($\LH = 0$ for quadratic $f$)
\end{itemizeM}
\item Damped NR with backtracking: $\alpha_{\text{init}} = 1$, $\gamma_1 \in (0, 1/2)$, $\tau \in (0, 1)$
\begin{framed}
\textbf{Theorem \textnormal{(Convergence of damped NR)}.}
With the scaled gradient norm $\theta^{[t]} = \frac{\LH}{2m^2}\|\nabla f(\xv^{[t]})\|$, there exist
$\eta \in (0, m^2/\LH]$ and $\delta > 0$ such that
\begin{align*}
\|\nabla f(\xv^{[t]})\| \ge \eta \; &\Rightarrow \; f(\xv^{[t+1]}) - f(\xv^{[t]}) \le - \delta
&& \textit{(damped)} \\
\|\nabla f(\xv^{[t]})\| < \eta \; &\Rightarrow \; \alpha^{[t]} = 1 \, \text{ and } \,
\theta^{[t+1]} \le \big( \theta^{[t]} \big)^2 && \textit{(pure)}
\end{align*}
\end{framed}
\end{framei}

\begin{framei}{Proof strategy}
\item tbd
\end{framei}

\endlecture
\end{document}
